\begin{lemma}
Let \(G\) be a finite graph, let \(\Delta\in\mathbb N\), and let
\(0<\gamma\le\Delta\).  There is a probability law \(\nu\) on pairs
\((A,\pi)\), where \(A\subseteq V(G)\) and \(\pi:V(G)\to\mathbb R\), such
that:
\[
\nu(A=A_0)=
\left(\frac{\gamma}{\Delta}\right)^{|A_0|}
\left(1-\frac{\gamma}{\Delta}\right)^{|V(G)|-|A_0|},
\]
\[
\nu(A=A_0,\ \pi(v)\in[0,1]\text{ for all }v\in A_0)=\nu(A=A_0),
\]
for every \(t\in[0,1]^{V(G)}\),
\[
\nu(A=A_0,\ 0\le \pi(v)\le t_v\text{ for all }v)
=\nu(A=A_0)\prod_{v\in V(G)}t_v,
\]
and, for every \(v\in A_0\),
\[
\nu\bigl(A=A_0,\ \pi(u)\in[0,1]\text{ for all }u\in A_0,\ 
\pi(u)<\pi(v)\text{ for all }u\in A_0\cap N_G(v)\bigr)
\]
equals
\[
\nu(A=A_0)\int_0^1 a^{|A_0\cap N_G(v)|}\,da .
\]
\end{lemma}
\begin{proof}
Put \(p=\gamma/\Delta\).  By
\noderef{SamplingBernoulliActivationMass}, the activation weights
\[
p^{|A|}(1-p)^{|V(G)|-|A|}
\]
are nonnegative and sum to one over all \(A\subseteq V(G)\).  Thus they define
a probability law for the activation set \(A\).

On the finite product cube \([0,1]^{V(G)}\), take the product of one
Lebesgue-uniform probability measure for each vertex coordinate.  Because the
indexing set \(V(G)\) is finite, this product probability measure exists and
has total mass one.  Its defining finite-dimensional rectangle formula gives,
for every vector \(t\in[0,1]^{V(G)}\), the mass
\[
\prod_{v\in V(G)} t_v
\]
of the lower rectangle \(0\le \pi(v)\le t_v\) for all \(v\).  It is supported
on the unit cube, so restricting the activated coordinates to \([0,1]\) loses
no mass.

Let \(\nu\) be the product of the activation law and this priority law.  The
total mass of \(\nu\) is \(1\cdot1=1\).  The formula for \(\nu(A=A_0)\) is the
activation atom weight, and the lower-rectangle identity follows by
multiplying that atom weight by the priority lower-rectangle mass.  This gives
the first three displayed identities.

It remains to verify the fixed-activation comparison identity.  Fix
\(A_0\subseteq V(G)\) and \(v\in A_0\).  Conditional on the activation atom
\(A=A_0\), the priority coordinates are independent uniforms on \([0,1]\).
Condition further on \(\pi(v)=a\).  For each activated neighbor
\(u\in A_0\cap N_G(v)\), the condition \(\pi(u)<\pi(v)\) then has probability
\(a\), and these conditions are independent over the distinct neighbors.
Coordinates not in \(A_0\cap N_G(v)\) impose no restriction.  Therefore the
conditional probability is \(a^{|A_0\cap N_G(v)|}\).  Integrating over the
uniform value \(a\in[0,1]\) of \(\pi(v)\), and then multiplying by the atom
mass \(\nu(A=A_0)\), gives exactly the displayed integral formula.
\end{proof}
